% Documentação técnica - Seguro Viagem (Gestão de Apólices)
% Compilar dentro de docs/latex:  pdflatex documentacao.tex  (duas vezes, para o sumário)
\documentclass[11pt,a4paper]{article}

\usepackage[T1]{fontenc}
\usepackage[utf8]{inputenc}
\usepackage[brazil]{babel}
\usepackage{lmodern}
\usepackage[a4paper,margin=2.5cm]{geometry}
\usepackage{listings}
\usepackage{booktabs}
\usepackage{tabularx}
\usepackage{amsmath}
\usepackage{tikz}
\usetikzlibrary{positioning,arrows.meta}
\usepackage[hidelinks]{hyperref}

\setlength{\parindent}{0pt}
\setlength{\parskip}{0.5em}
\setlength{\emergencystretch}{3em}

\newcommand{\arq}[1]{\texttt{\detokenize{#1}}}

\lstset{
  basicstyle=\ttfamily\small,
  keywordstyle=\bfseries,
  commentstyle=\itshape,
  numbers=left, numberstyle=\scriptsize, numbersep=6pt,
  frame=lines, xleftmargin=16pt,
  breaklines=true, columns=fullflexible, keepspaces=true,
  showstringspaces=false, tabsize=4,
  inputencoding=utf8, extendedchars=true,
  literate=
    {á}{{\'a}}1 {é}{{\'e}}1 {í}{{\'i}}1 {ó}{{\'o}}1 {ú}{{\'u}}1
    {Á}{{\'A}}1 {É}{{\'E}}1 {â}{{\^a}}1 {ê}{{\^e}}1 {ô}{{\^o}}1
    {ã}{{\~a}}1 {õ}{{\~o}}1 {ç}{{\c{c}}}1 {—}{{---}}1,
}
\lstdefinestyle{php}{language=PHP, morekeywords={enum,match,readonly,fn,case,interface,implements}}
\lstdefinestyle{txt}{numbers=none}
\renewcommand{\lstlistingname}{Código}

\newcommand{\codigo}[2]{\lstinputlisting[style=php,caption={\arq{#1}}]{../../#2}}

\title{Seguro Viagem -- Gestão de Apólices\\[0.3em]
  \large Documentação técnica do código}
\author{Carlos Guilherme Fontes Pereira\\ \small Teste técnico -- CORIS}
\date{\today}

\begin{document}
\maketitle

\begin{abstract}
Este documento descreve o funcionamento de uma aplicação web para gestão de apólices de seguro
viagem, construída com Laravel 12 (PHP 8.3), React 18 e MySQL 8, com implantação prevista na
Microsoft Azure. São apresentados o modelo de dados, as regras de negócio, a organização do
código em camadas, a API e os testes automatizados, com o objetivo de permitir que o leitor
entenda o sistema e consiga modificá-lo.
\end{abstract}

\tableofcontents
\newpage

% =====================================================================
\section{Introdução}

A aplicação permite a um operador de seguradora \textbf{cadastrar, consultar, editar e excluir}
apólices de seguro viagem. Ao preencher o formulário, o operador informa os dados do segurado,
o destino, o plano e o período da viagem; o sistema então calcula o \emph{prêmio}, isto é, o
valor a ser pago pelo seguro.

O projeto está dividido em duas partes independentes que conversam por HTTP/JSON:

\begin{itemize}
  \item \textbf{backend} (\arq{backend/}): API REST em Laravel, responsável pelas regras e pelo banco;
  \item \textbf{frontend} (\arq{frontend/}): interface em React, que apenas exibe dados e envia formulários.
\end{itemize}

\subsection{Tecnologias}

\begin{center}
\begin{tabular}{@{}ll@{}}
\toprule
Camada & Tecnologia \\
\midrule
Backend & PHP 8.3, Laravel 12 \\
Frontend & React 18, Vite, React Router \\
Banco de dados & MySQL 8 (SQLite no desenvolvimento local) \\
Testes & PHPUnit \\
Ambiente & Docker Compose \\
Nuvem & Microsoft Azure \\
\bottomrule
\end{tabular}
\end{center}

% =====================================================================
\section{Modelo de dados}

O banco possui duas tabelas com relacionamento \textbf{1:N}: um segurado pode ter várias
apólices, e cada apólice pertence a um único segurado.

\begin{center}
\begin{tikzpicture}[tab/.style={draw, align=left, font=\ttfamily\small, inner sep=5pt}]
  \node[tab] (s) {\textbf{segurados}\\[2pt]
    id (PK)\\ nome\\ cpf (único)\\ email\\ data\_nascimento};
  \node[tab, right=3cm of s] (a) {\textbf{apolices}\\[2pt]
    id (PK)\\ numero (único)\\ segurado\_id (FK)\\ destino, plano\\
    inicio\_vigencia, fim\_vigencia\\ valor\_premio\_centavos\\ status\\ deleted\_at};
  \draw (s.east) -- node[above, font=\small]{1 \hspace{2cm} N} (s.east -| a.west);
\end{tikzpicture}
\end{center}

Três decisões merecem destaque:

\begin{enumerate}
  \item \textbf{Valores em centavos.} O prêmio é guardado como número inteiro de centavos
        (\texttt{32370} = R\$~323,70). Números de ponto flutuante não representam decimais com
        exatidão (em \texttt{float}, $0{,}1 + 0{,}2 = 0{,}30000000000000004$), o que é inaceitável
        para valores financeiros.
  \item \textbf{Segurado separado da apólice.} O CPF é único, então uma mesma pessoa que contrata
        outro seguro é reaproveitada, sem duplicar dados.
  \item \textbf{Exclusão lógica.} Excluir uma apólice apenas preenche a coluna \texttt{deleted\_at}.
        O registro some das consultas, mas continua no banco, pois contratos de seguro precisam ser
        mantidos para auditoria.
\end{enumerate}

As tabelas são criadas pelas \emph{migrations} em \arq{backend/database/migrations/}.

% =====================================================================
\section{Regras de negócio}

\subsection{Cálculo do prêmio}

\begin{equation*}
  \text{prêmio} = \text{diária do plano} \times \text{dias}
  \times \text{fator do destino} \times \text{fator da idade}
\end{equation*}

\begin{center}
\begin{tabular}{@{}lr@{}}
\toprule
Plano & Diária \\
\midrule
Essencial & R\$ 12,90 \\
Plus & R\$ 24,90 \\
Premium & R\$ 39,90 \\
\bottomrule
\end{tabular}
\quad
\begin{tabular}{@{}lr@{}}
\toprule
Destino & Fator \\
\midrule
Nacional & 50\% \\
América do Sul & 100\% \\
Europa & 130\% \\
América do Norte & 140\% \\
Ásia, África, Oceania & 150\% \\
\bottomrule
\end{tabular}
\quad
\begin{tabular}{@{}lr@{}}
\toprule
Idade & Fator \\
\midrule
até 59 & 100\% \\
60 a 74 & 160\% \\
75 ou mais & 250\% \\
\bottomrule
\end{tabular}
\end{center}

Os dias incluem o primeiro e o último (de 01/10 a 10/10 são 10 dias), e a idade é calculada na
data de início da viagem.

\textbf{Exemplo.} Plano Plus, Europa, 10 dias, segurado de 31 anos:
\[
  2490 \times 10 = 24900 \;\rightarrow\; 24900 \times 1{,}30 = 32370 \;\rightarrow\; 32370 \times 1{,}00 = 32370
  \text{ centavos} = \text{R\$ 323,70}.
\]

\subsection{Validações e demais regras}

\begin{itemize}
  \item O CPF deve ter dígitos verificadores válidos.
  \item O fim da vigência não pode ser anterior ao início, e a viagem tem no máximo 365 dias.
  \item O início da vigência não pode ser anterior a hoje. Na edição, a regra só é aplicada se a
        data de início mudou, para que uma apólice já em vigor continue editável.
  \item Toda apólice recebe um número único no formato \texttt{CRS-2026-XXXXXXXX}.
  \item Uma apólice cancelada só pode ser alterada para voltar a ficar ativa.
\end{itemize}

% =====================================================================
\section{Arquitetura do backend}

O backend é organizado em camadas, e cada classe tem uma única responsabilidade:

\begin{center}
\begin{tikzpicture}[
  c/.style={draw, minimum width=2.4cm, minimum height=0.8cm, font=\small},
  s/.style={-{Stealth}}, node distance=0.6cm]
  \node[c] (r) {Rota};
  \node[c, right=of r] (f) {Form Request};
  \node[c, right=of f] (co) {Controller};
  \node[c, right=of co] (se) {Service};
  \node[c, below=0.7cm of se] (ca) {Calculadora};
  \node[c, left=of ca] (m) {Model};
  \node[c, left=of m] (re) {Resource};
  \draw[s] (r) -- (f); \draw[s] (f) -- (co); \draw[s] (co) -- (se);
  \draw[s] (se) -- (ca); \draw[s] (se) -- (m); \draw[s] (co) -- (re);
\end{tikzpicture}
\end{center}

\begin{center}
\begin{tabularx}{\linewidth}{@{}lX@{}}
\toprule
Camada & Responsabilidade \\
\midrule
Rota (\arq{routes/api.php}) & associa a URL ao método do controller \\
Form Request (\arq{app/Http/Requests}) & valida a entrada; se falhar, responde 422 \\
Controller (\arq{app/Http/Controllers/Api}) & recebe a requisição e devolve a resposta HTTP \\
Service (\arq{app/Services/ApoliceService.php}) & aplica as regras de negócio \\
Calculadora (\arq{app/Services/Premio}) & calcula o prêmio \\
Model (\arq{app/Models}) & representa as tabelas (Eloquent) \\
Resource (\arq{app/Http/Resources}) & define o formato do JSON de saída \\
Enums (\arq{app/Enums}) & planos, destinos e status, com seus valores \\
\bottomrule
\end{tabularx}
\end{center}

\subsection{Exemplo: o que acontece ao emitir uma apólice}

\begin{enumerate}
  \item O React envia \texttt{POST /api/apolices} com os dados do formulário em JSON.
  \item O \arq{SalvarApoliceRequest} valida os campos. Havendo erro, a API responde \texttt{422}
        com uma mensagem por campo, e o fluxo termina aqui.
  \item O \arq{ApoliceController} chama \texttt{ApoliceService::criar}.
  \item O service, dentro de uma transação, cria ou atualiza o segurado pelo CPF, calcula o
        prêmio, gera o número da apólice e a grava.
  \item O \arq{ApoliceResource} converte a apólice em JSON, e a resposta é \texttt{201 Created}.
\end{enumerate}

\subsection{Rotas}

\codigo{routes/api.php}{backend/routes/api.php}

As rotas \texttt{resumo} e \texttt{cotacao} vêm antes do \texttt{apiResource}, e
\texttt{whereNumber} impede que \texttt{/apolices/resumo} seja interpretado como
\texttt{/apolices/\{id\}}.

\subsection{Validação}

\arq{CotacaoRequest} contém as regras comuns à cotação e ao cadastro. \arq{SalvarApoliceRequest}
herda essas regras e acrescenta o campo \texttt{status} e a verificação da data de início.
Validações que envolvem mais de um campo ficam no método \texttt{after()}.

\codigo{app/Http/Requests/CotacaoRequest.php}{backend/app/Http/Requests/CotacaoRequest.php}

O CPF é validado pela regra \arq{app/Rules/Cpf.php}: cada dígito verificador é obtido
multiplicando os dígitos anteriores por pesos decrescentes, somando e calculando
$(10 \cdot \text{soma}) \bmod 11$ (o resultado 10 vale 0).

\subsection{Controller}

O controller não contém regra de negócio: apenas repassa os dados validados ao service e escolhe
o código HTTP da resposta.

\codigo{app/Http/Controllers/Api/ApoliceController.php}{backend/app/Http/Controllers/Api/ApoliceController.php}

\subsection{Service}

É a classe central do sistema. Observe que o construtor recebe a \emph{interface}
\texttt{CalculadoraPremio}, e não uma classe concreta.

\codigo{app/Services/ApoliceService.php}{backend/app/Services/ApoliceService.php}

Pontos importantes:
\begin{itemize}
  \item \texttt{DB::transaction} garante que segurado e apólice sejam gravados juntos, ou nenhum dos dois.
  \item \texttt{updateOrCreate} reaproveita o segurado quando o CPF já existe.
  \item \texttt{gerarNumero} consulta também as apólices excluídas (\texttt{withTrashed}), para nunca repetir um número.
\end{itemize}

\subsection{Calculadora do prêmio}

\codigo{app/Services/Premio/CalculadoraPremio.php}{backend/app/Services/Premio/CalculadoraPremio.php}
\codigo{app/Services/Premio/CalculadoraPremioViagem.php}{backend/app/Services/Premio/CalculadoraPremioViagem.php}

O arredondamento usa apenas inteiros: \texttt{intdiv(centavos * percentual + 50, 100)}. Somar 50
antes da divisão inteira arredonda para o centavo mais próximo.

A ligação entre a interface e a implementação é feita em
\arq{app/Providers/AppServiceProvider.php}:

\begin{lstlisting}[style=php,numbers=none]
$this->app->bind(CalculadoraPremio::class, CalculadoraPremioViagem::class);
\end{lstlisting}

\subsection{Models e Enums}

\arq{Segurado} possui \texttt{hasMany(Apolice)} e \arq{Apolice} possui \texttt{belongsTo(Segurado)}.
A apólice usa \texttt{SoftDeletes} (exclusão lógica) e \emph{casts}, que convertem automaticamente
as colunas em enums e datas. O \emph{scope} \texttt{buscar} procura um termo no número da apólice e
no nome, e-mail ou CPF do segurado.

Os valores de negócio ficam nos enums. Por exemplo, a diária de cada plano:

\begin{lstlisting}[style=php,numbers=none]
public function diariaCentavos(): int
{
    return match ($this) {
        self::Essencial => 1290,
        self::Plus => 2490,
        self::Premium => 3990,
    };
}
\end{lstlisting}

Para criar um novo plano basta adicionar um \texttt{case}: o frontend obtém a lista pela rota
\texttt{/api/opcoes} e passa a exibi-lo sem alteração.

% =====================================================================
\section{SOLID e Clean Code}

\begin{center}
\begin{tabularx}{\linewidth}{@{}lX@{}}
\toprule
Princípio & Aplicação \\
\midrule
S -- Responsabilidade única & validação, HTTP, regra de negócio, cálculo e formatação em classes distintas \\
O -- Aberto/fechado & nova regra de preço é uma nova classe que implementa \texttt{CalculadoraPremio}; novo plano é um novo \texttt{case} \\
L -- Substituição de Liskov & qualquer implementação da calculadora pode substituir a atual \\
I -- Segregação de interfaces & a interface da calculadora possui um único método \\
D -- Inversão de dependência & o service depende da interface; o provider escolhe a implementação \\
\bottomrule
\end{tabularx}
\end{center}

Quanto a Clean Code: os nomes seguem o vocabulário do negócio (apólice, segurado, prêmio,
vigência), os métodos são curtos e os valores fixos ficam em enums ou constantes
(\texttt{VIGENCIA\_MAXIMA\_DIAS}, \texttt{POR\_PAGINA}).

% =====================================================================
\section{API}

\begin{center}
\begin{tabularx}{\linewidth}{@{}llX@{}}
\toprule
Método & Rota & Descrição \\
\midrule
GET & \texttt{/api/apolices} & lista paginada (\texttt{busca}, \texttt{status}, \texttt{page}) \\
GET & \texttt{/api/apolices/resumo} & totais da tela inicial \\
GET & \texttt{/api/apolices/\{id\}} & detalhe \\
POST & \texttt{/api/apolices} & cadastro (201) \\
PUT & \texttt{/api/apolices/\{id\}} & edição, com recálculo do prêmio \\
DELETE & \texttt{/api/apolices/\{id\}} & exclusão lógica (204) \\
POST & \texttt{/api/apolices/cotacao} & calcula o prêmio sem salvar \\
GET & \texttt{/api/opcoes} & planos, destinos e status \\
\bottomrule
\end{tabularx}
\end{center}

Exemplo de requisição e resposta de erro:

\begin{lstlisting}[style=txt]
POST /api/apolices
{ "seguradoNome": "Carlos Pereira", "seguradoCpf": "123.456.789-00",
  "seguradoEmail": "carlos@email.com", "seguradoNascimento": "1995-05-10",
  "destino": "europa", "plano": "ouro",
  "inicioVigencia": "2026-10-01", "fimVigencia": "2026-10-10" }

422 Unprocessable Content
{ "message": "CPF inválido. (e mais 1 erro)",
  "errors": { "seguradoCpf": ["CPF inválido."], "plano": ["Plano inválido."] } }
\end{lstlisting}

O banco usa nomes em \texttt{snake\_case} e a API responde em \texttt{camelCase}; essa conversão
é feita pelo \arq{ApoliceResource}.

% =====================================================================
\section{Frontend}

\begin{center}
\begin{tabularx}{\linewidth}{@{}lX@{}}
\toprule
Pasta & Conteúdo \\
\midrule
\arq{src/pages} & telas: lista, detalhe, cadastro e edição \\
\arq{src/components} & formulário, paginação, diálogo de confirmação, notificações \\
\arq{src/api} & \texttt{client.js} (função \texttt{request}) e \texttt{apolices.js} (uma função por endpoint) \\
\arq{src/hooks} & \texttt{useApolice} e \texttt{useOpcoes} (busca de dados) \\
\arq{src/utils/format.js} & formatação de moeda, datas e máscara de CPF \\
\bottomrule
\end{tabularx}
\end{center}

Toda chamada à API passa por \texttt{request}, que define os cabeçalhos, converte o corpo para
JSON e, em caso de erro, lança um \texttt{ApiError} com a mensagem de cada campo, exibida abaixo do
campo correspondente no formulário.

A \textbf{cotação em tempo real} funciona assim: quando todos os campos estão preenchidos, o
formulário aguarda 400\,ms sem alterações (\emph{debounce}) e chama \texttt{/api/apolices/cotacao}.
Se o usuário continuar digitando, a requisição anterior é cancelada com \texttt{AbortController},
evitando que uma resposta antiga substitua a mais recente. A busca da lista usa a mesma técnica.

% =====================================================================
\section{Testes}

São 25 testes automatizados (PHPUnit), executados com \texttt{php artisan test} na pasta
\arq{backend}.

\begin{itemize}
  \item \textbf{Unitários:} \arq{CalculadoraPremioViagemTest} verifica cinco combinações de plano,
        destino e idade; \arq{CpfTest} verifica CPFs válidos e inválidos.
  \item \textbf{De API:} \arq{ApoliceApiTest} faz requisições reais e verifica cadastro, reaproveitamento
        do segurado, busca e paginação, edição com recálculo, regra da apólice cancelada, exclusão
        lógica, mensagens de validação, limites da vigência, cotação e resumo.
\end{itemize}

Nos testes de API a data atual é fixada com \texttt{Carbon::setTestNow('2026-09-20')}, para que o
resultado não dependa do dia em que o teste é executado, e o banco é recriado a cada teste
(\texttt{RefreshDatabase}).

% =====================================================================
\section{Execução e implantação}

\subsection{Com Docker}

\begin{lstlisting}[style=txt]
docker compose up -d --build
docker compose exec api php artisan db:seed     (dados de exemplo)
\end{lstlisting}

Aplicação em \texttt{http://localhost:3000} e API em \texttt{http://localhost:8000/api}.

\subsection{Na Azure}

\begin{center}
\begin{tikzpicture}[
  c/.style={draw, align=center, font=\small, minimum width=3.6cm, minimum height=0.9cm},
  s/.style={-{Stealth}}, node distance=0.9cm]
  \node[c] (u) {Usuário\\ (navegador)};
  \node[c, right=1.4cm of u] (swa) {Static Web Apps\\ (React)};
  \node[c, below=of swa] (app) {App Service\\ (API Laravel)};
  \node[c, below=of app] (db) {Database for MySQL};
  \node[c, right=1.4cm of app] (kv) {Key Vault\\ (segredos)};
  \draw[s] (u) -- (swa);
  \draw[s] (u) |- (app);
  \draw[s] (app) -- (db);
  \draw[s] (app) -- (kv);
\end{tikzpicture}
\end{center}

\begin{itemize}
  \item \textbf{Static Web Apps} hospeda o React compilado.
  \item \textbf{App Service} executa a API.
  \item \textbf{Azure Database for MySQL} armazena os dados.
  \item \textbf{Key Vault} guarda a senha do banco e a \texttt{APP\_KEY}, lidas pelo App Service via
        \emph{Managed Identity}.
\end{itemize}

% =====================================================================
\section{Trabalhos futuros}

\begin{itemize}
  \item Autenticação e perfis de acesso (Laravel Sanctum).
  \item Endosso: histórico das alterações de cada apólice.
  \item Infraestrutura como código (Bicep).
  \item Testes automatizados no frontend.
\end{itemize}

\end{document}
